\section{Discussion}
\label{sec:discussion}

\para{Verdict on the hypothesis.} The go-to-sleep mechanism is \emph{partially unified, with
a clear hierarchy of drivers: model $\gg$ situation $>$ person}. The situation is the
dominant within-model driver (H1, strongly supported; $\omega^2 = 0.11$, P from $0.42$ to
$0.91$). A person effect exists for occupation, vulnerability, and age but not gender (H2,
supported but small; $V \approx 0.14$--$0.19$), and it is about $3\times$ smaller than the
situation effect and far smaller than the model effect. The scenario~$\times$~persona
interaction is significant and itself exceeds the main person effect ($\omega^2 = 0.042$;
H4 supported), meaning the same identity cue matters more in some scenarios than others.
The cross-model dimension is where the behavior is \emph{least} unified (H5): a $0.15$ to
$0.82$ spread in level and only moderate agreement on the conditioning pattern.

\para{The surprise: physiological salience, not vulnerability.} The most novel result is the
refutation of H3, the protective-stereotype prediction drawn from the literature. The person
effect is not organized by demographic vulnerability; it tracks \emph{bodily stakes and the
legitimacy of being awake}. Athletes, the ill, and children --- all of whom supply a salient
physiological reason to rest --- get more sleep advice, while the night-shift nurse, whose
wakefulness is part of her job, gets less. This is closer to the role-dependent finding of
\citet{goalpref2026} than to flat demographic bias: the model appears to reason about
\emph{why} the user is awake and whether sleeping now is locally sensible, rather than
applying a category-level ``protect the vulnerable'' rule.

\para{Relation to prior work.} Our results are consistent with \citet{eloundou2024fairness}
(the person channel is real but modest, and gender effects are weak or absent here) and with
\citet{differentcues2026} (single-cue, stereotype-shaped conclusions are fragile --- our
naïve stereotype contrast failed and reversed). We extend the personalization-bias
line~\citep{vijjini2024safety} by showing that, for protective wellbeing advice, the
operative variable is situational legitimacy of wakefulness rather than demographic category.
The dominant model effect echoes a recurring lesson: claims about ``what LLMs do'' rarely
generalize across providers, and the biggest predictor of whether a user is told to go to
bed is simply which model they are talking to.

\para{Robustness.} Refusals were negligible ($1.9\%$). Every headline conclusion ---
situation $>$ person, model dominance, and the non-stereotyped person effect --- holds across
all four outcome definitions and in the ordinal analysis, so none is an artifact of the
strength~$\geq 2$ threshold. The strict \texttt{insist\_sleep} and lexical \texttt{regex}
definitions, which favor no particular panel member, give the same model ordering, which
mitigates the self-preference concern of using an OpenAI judge alongside an OpenAI model.

\para{Limitations.} (1) \emph{Single judge.} One judge model graded all responses; manual
validation found it reliable at strength $0$/$3$ but noisier at $1{\leftrightarrow}2$. We
mitigate with a strength~$\geq 2$ primary outcome and four-way sensitivity, but a multi-judge
panel with inter-annotator $\kappa$ would strengthen the labels. (2) \emph{Explicit, single,
English cues.} Identity is stated explicitly (``I'm 9 years old''); implicit or name-based
cues, known to yield smaller effects, and non-English prompts were not tested, so the
reported person effect is likely an upper bound on real-world implicit conditioning. (3)
\emph{Synthetic probe.} Eight templated scenarios approximate but do not equal real extended
sessions with multi-turn fatigue build-up. (4) \emph{Model snapshot.} Provider versions
rotate; these results are a June-2026 snapshot of pinned IDs. (5) \emph{Scope of person.} We
covered age, gender, occupation, and vulnerability; race, religion, politics, and disability
were out of scope.

\para{Broader implications.} ``LLMs tell people to go to sleep'' is not a property of LLMs
but of a \emph{particular} model's late-night disposition. The user-facing fairness concern
--- differential protective advice by identity --- is real but modest, and is organized by
physiological salience rather than demographic vulnerability. For deployers, the practical
takeaway is that wellbeing nudges are not a uniform safety reflex: their presence, strength,
and conditioning vary by an order of magnitude across models, and any policy that assumes a
consistent ``go to sleep'' behavior across a model fleet is mistaken.
